\documentclass[11pt]{article}
\usepackage[utf8]{inputenc}
\usepackage[T1]{fontenc}
\usepackage{lmodern}
\usepackage[margin=1.05in]{geometry}
\usepackage{amsmath,amssymb,amsthm,mathtools}
\usepackage{microtype,booktabs,array,longtable}
\usepackage{listings}
\usepackage{xcolor}
\usepackage[colorlinks=true,linkcolor=blue!45!black,citecolor=blue!45!black,urlcolor=blue!45!black]{hyperref}
\usepackage[nameinlink,capitalise,noabbrev]{cleveref}
\numberwithin{equation}{section}
\newtheorem{theorem}{Theorem}[section]
\newtheorem{proposition}[theorem]{Proposition}
\newtheorem{lemma}[theorem]{Lemma}
\newtheorem{corollary}[theorem]{Corollary}
\theoremstyle{definition}
\newtheorem{definition}[theorem]{Definition}
\theoremstyle{remark}
\newtheorem{remark}[theorem]{Remark}
\newcommand{\F}{\mathbb F}
\newcommand{\PG}{\mathrm{PG}}
\newcommand{\PGL}{\mathrm{PGL}}
\newcommand{\GL}{\mathrm{GL}}
\newcommand{\wt}{\mathrm{wt}}
\lstset{basicstyle=\ttfamily\small,columns=fullflexible,keepspaces=true,breaklines=true,
  frame=single,framerule=0.3pt,xleftmargin=4pt,xrightmargin=4pt,
  literate={⟨}{{$\langle$}}1 {⟩}{{$\rangle$}}1 {→}{{$\to$}}1 {≠}{{$\neq$}}1 {≤}{{$\le$}}1
           {∀}{{$\forall$}}1 {∧}{{$\wedge$}}1 {∃}{{$\exists$}}1 {↔}{{$\leftrightarrow$}}1
           {×}{{$\times$}}1 {…}{{\ldots}}1 {·}{{$\cdot$}}1 {₃}{{$_3$}}1 {²}{{$^2$}}1}
\title{New four-dimensional linear codes over $\mathbb{F}_8$ and $\mathbb{F}_9$ improving Maruta's tables}
\author{Claude (AI model by Anthropic)\\ \small run and submitted by Aleksandr Bersenev}
\date{October 8, 2026\\[3pt]\footnotesize Produced autonomously by Claude (Opus 5.5) agents in Claude Code; run and submitted by Aleksandr Bersenev.\\ \footnotesize Not an official Anthropic publication. Not peer reviewed.}
\hypersetup{pdftitle={New four-dimensional linear codes over GF(8) and GF(9) improving Maruta's tables},pdfauthor={Claude (AI model by Anthropic); submitted by Aleksandr Bersenev}}
\begin{document}
\maketitle
\begin{abstract}
Let $n_q(4,d)$ be the smallest length of a linear $[n,4,d]_q$ code. We give 32 explicit
four-dimensional linear codes, 24 over $\F_9$ and 8 over $\F_8$. Together with the codes obtained
from them by puncturing, they are shorter than the upper bound recorded in Maruta's online tables
of $n_q(4,d)$ (as of October 8, 2026) for $191$ values of $d$: for $149$ values of $d$ with
$190\le d\le 513$ when $q=9$ (twenty of them by $2$), and for $42$ values of $d$ with
$137\le d\le 376$ when $q=8$. In $29$ cases the new length equals the lower bound proved in the
literature, which gives
$n_8(4,d)=g_8(4,d)+1$ for $d\in\{337,\dots,345\}\cup\{353,354\}\cup\{361,\dots,376\}$ and
$n_9(4,512)=578$, $n_9(4,513)=579$, where $g_q(4,d)$ is the Griesmer bound.
The codes were found by a prescribed-automorphism search (orbits of cyclic subgroups of
$\PGL(4,q)$, integer programming, and greedy lengthening). The existence of every code is certified
in Lean~4 by a kernel-checked computation whose soundness is proved in Lean. The lower bounds are
cited, not re-proved.
\end{abstract}
\tableofcontents

\section{Introduction}

For a prime power $q$ and integers $k,d\ge 1$, let $n_q(k,d)$ denote the minimum length $n$ for
which a linear $[n,k,d]_q$ code exists, i.e.\ a $k$-dimensional subspace of $\F_q^n$ all of whose
nonzero vectors have Hamming weight at least $d$. Determining $n_q(k,d)$ is one of the basic
problems of coding theory (see the surveys of Hill~\cite{Hill1992}). The classical lower bound is
the Griesmer bound~\cite{Griesmer1960,SolomonStiffler1965}
\[
  n_q(k,d)\;\ge\; g_q(k,d)\;:=\;\sum_{i=0}^{k-1}\Big\lceil \frac{d}{q^{i}}\Big\rceil .
\]
For $k\le 3$ and small $q$ the function $n_q(k,d)$ is known. For $k=4$ it is known for
$q\le 4$, but for $q\ge 5$ many values remain open.
T.~Maruta maintains online tables~\cite{MarutaTable} with the current
intervals $n_{\mathrm{lo}}(d)\le n_q(4,d)\le n_{\mathrm{up}}(d)$ for $q\in\{5,7,8,9\}$ and all $d$
below the range where $n_q(4,d)=g_q(4,d)$ is known.
\begin{itemize}
\item For $q=8$ the table follows Kanazawa--Maruta~\cite{KanazawaMaruta2011}, with later updates
(e.g.~\cite{BFM2021,MaLuo2024}) for values of $d$ not considered here. That paper determines
$n_8(4,d)$ for 117 values of $d$ with $113\le d\le 832$ and gives bounds for the others. In
particular \cite[Thm.~1.1(4)]{KanazawaMaruta2011} shows
$g_8(4,d)+1\le n_8(4,d)\le g_8(4,d)+2$ for $337\le d\le 376$.
\item For $q=9$ the table follows Kumegawa--Okazaki--Maruta~\cite{KOM2017} (a few lower bounds have been
updated since), who constructed many
mid-range codes ``with the aid of a computer, but the upper bounds are still weak to be improved''
\cite[p.~10]{KOM2017}. In particular \cite[Thm.~1.7(4)]{KOM2017} gives
$g_9(4,d)+1\le n_9(4,d)\le g_9(4,d)+2$ for $d=512,513$.
\end{itemize}
In the mid range the intervals are typically $g_9(4,d)\le n_9(4,d)\le g_9(4,d)+3$ or $+4$ for $q=9$, and
$g_8(4,d)+1\le n_8(4,d)\le g_8(4,d)+2$ for $q=8$. These lengths ($n\ge 160$) lie outside Grassl's
code tables~\cite{Grassl}, which stop at $n=130$ for $q=8,9$. They also lie outside MinT~\cite{MinT},
whose linear-code data are limited to co-dimension at most $173$ for base $8$ and $150$ for base $9$.

\paragraph{What is new.} We construct 32 linear codes with $k=4$ over $\F_8$ and $\F_9$. Together with
their punctured codes, they lower the recorded upper bound $n_{\mathrm{up}}(d)$ for 191 table entries
(\cref{thm:q9,thm:q8}). Twenty-nine of the new lengths coincide with lower bounds proved in
\cite{KanazawaMaruta2011,KOM2017}, which determines these values of $n_q(4,d)$ (Corollary~\ref{cor:exact}).
To our knowledge, as of October 8, 2026, none of these codes or parameters appeared before. We
compared against Maruta's table, the two papers above, Grassl's tables and MinT, and searched the
recent literature on $n_8(4,d)$ and $n_9(4,d)$~\cite{BFM2021,MaLuo2024,KKM2026}. For every improved
entry with $q=8$ the table row (lower and upper bound) is the one of Table~3 of~\cite{KanazawaMaruta2011},
and for every improved entry with $q=9$ the upper bound is the one of Table~2 of~\cite{KOM2017}; the
updates of the table since then concern other values of $d$ or only lower bounds. The most recent
nonexistence results~\cite[Thm.~1.3]{KKM2026} ($n_8(4,d)=g_8(4,d)+1$ for $d\in\{278,398,406\}$,
$n_9(4,d)=g_9(4,d)+1$ for $d\in\{583,808\}$, and $n_9(4,196)\ge g_9(4,196)+1=223$) are not yet
reflected in the table; they do not affect the entries below, except that the lower bound for
$d=196$, $q=9$, should now be read as $223$ instead of $222$. We cannot rule out unpublished
constructions.

\paragraph{Method in one paragraph.} By the standard correspondence~\cite{Hamada1993}, an
$[n,4,d]_q$ code is a multiset of $n$ points of $\PG(3,q)$ meeting every plane in at most $n-d$
points. Following the prescribed-automorphism approach of Braun--Kohnert--Wassermann~\cite{BKW2005}
and the ``orbits of projectivities'' of Maruta--Shinohara--Takenaka~\cite{MST2008}, we chose random
projectivities $T$ of moderate order. We computed the orbits of $\langle T\rangle$ on points and on
planes, and solved with HiGHS~\cite{HiGHS} the integer program ``choose a union of point orbits of
maximum size subject to all plane intersections being $\le r$''. We then lengthened the solution
greedily by single points. Every hit was checked by exhaustive enumeration of all $q^4-1$ nonzero
codewords, and all base codes were then certified in Lean~4~\cite{Lean4}.

\section{Definitions and main results}

\begin{definition}
Let $L=(c_1,\dots,c_n)$ be a list of vectors of $\F_q^4$ (the columns of a $4\times n$ matrix $G$).
For $u\in\F_q^4$ the codeword of $u$ is $uG=(u\cdot c_1,\dots,u\cdot c_n)$. We say that $L$
\emph{generates an $[n,4,\ge d]_q$ code} if $\wt(uG)\ge d$ for all $u\ne 0$. Since $d\ge1$, the map
$u\mapsto uG$ is then injective, so the row space of $G$ is a linear code of length $n$, dimension
exactly $4$ and minimum distance at least $d$. Columns may repeat (non-projective codes), as in the
definition of $n_q(k,d)$.
\end{definition}

\begin{lemma}[puncturing]\label{lem:punct}
If $L$ generates an $[n,4,\ge d]_q$ code and $0\le t<d$, then deleting any $t$ columns gives an
$[n-t,4,\ge d-t]_q$ code. Hence $n_q(4,d-t)\le n_q(4,d)-t$.
\end{lemma}
\begin{proof}
Deleting $t$ coordinates lowers every weight by at most $t$. Since $d-t\ge1$, injectivity is preserved.
\end{proof}

\begin{theorem}[$q=9$]\label{thm:q9}
For each of the $149$ values of $d$ in \cref{tab:q9}, there is a linear $[n,4,\ge d]_9$ code whose
length $n$ (column ``new'') is strictly smaller than the upper bound $n_{\mathrm{up}}(d)$ recorded
in~\cite{MarutaTable}. The improved values of $d$ form the ranges
$190$--$192$, $194$--$200$, $231$--$232$, $241$--$243$, $253$--$256$, $262$--$279$, $289$--$292$,
$298$--$312$, $333$--$360$, $370$--$375$, $379$--$387$, $397$--$417$, $424$--$429$, $433$--$438$,
$487$--$489$, $496$--$500$, $505$--$513$. For twenty of them the improvement is by $2$, for example
$[312,4,274]_9$ (previous upper bound $314$), $[345,4,303]_9$ ($347$), $[380,4,334]_9$ ($382$),
$[390,4,343]_9$ ($392$), $[420,4,370]_9$ ($422$), $[460,4,406]_9$ ($462$), $[468,4,414]_9$ ($470$).
\end{theorem}

\begin{theorem}[$q=8$]\label{thm:q8}
For each of the $42$ values of $d$ in \cref{tab:q8}, namely $137$--$140$, $221$--$224$,
$329$--$345$, $353$--$354$ and $361$--$376$, there is a linear $[n,4,\ge d]_8$ code with $n$
strictly smaller than $n_{\mathrm{up}}(d)$ in~\cite{MarutaTable}.
\end{theorem}

\begin{corollary}\label{cor:exact}
Combining \cref{thm:q8,thm:q9} with the lower bounds of \cite[Thm.~1.1(4)]{KanazawaMaruta2011}
and \cite[Thm.~1.7(4)]{KOM2017} gives
\[
  n_8(4,d)=g_8(4,d)+1 \quad\text{for } d\in\{337,\dots,345\}\cup\{353,354\}\cup\{361,\dots,376\},
\]
and $n_9(4,512)=578=g_9(4,512)+1$, $n_9(4,513)=579=g_9(4,513)+1$.
\end{corollary}
\begin{proof}
By the cited theorems, $n\ge g+1$ for these $d$. For the upper bound we have explicit codes:
\begin{itemize}
\item $[395,4,344]_8$ and its punctured codes $[395-t,4,344-t]_8$ for $t\le 7$ cover $d=337,\dots,344$;
\item $[397,4,345]_8$ covers $d=345$;
\item $[407,4,354]_8$ and its punctured code $[406,4,353]_8$ cover $d=353,354$;
\item $[422,4,368]_8$ and its punctured codes cover $d=361,\dots,368$;
\item $[431,4,376]_8$ and its punctured codes cover $d=369,\dots,376$;
\item $[579,4,513]_9$ and $[578,4,512]_9$ cover the two values for $q=9$.
\end{itemize}
In each case the length equals $g_q(4,d)+1$, which a direct computation confirms; for example
$g_8(4,344)=344+43+6+1=394$.
\end{proof}

The lower bounds used in Corollary~\ref{cor:exact} are \emph{not} re-proved here; the Lean development
certifies only the existence of the codes (\cref{sec:lean}).

\section{The construction}\label{sec:construction}

\subsection{Geometric reformulation}
Identify a nonzero column $c\in\F_q^4$ with the point $\langle c\rangle$ of $\PG(3,q)$ and a nonzero
message $u$ with the plane $u^\perp=\{\langle x\rangle: u\cdot x=0\}$. If $M$ is the multiset of
columns, then $\wt(uG)=n-|M\cap u^\perp|$. So $L$ generates an $[n,4,\ge n-r]_q$ code iff every
plane contains at most $r$ points of $M$, counted with multiplicity. $\PG(3,q)$ has
$\theta=q^3+q^2+q+1$ points and planes, i.e.\ $585$ for $q=8$ and $820$ for $q=9$.

\subsection{Prescribed automorphisms and an integer program}
Let $T\in\GL(4,q)$, and let $\langle T\rangle$ act on points by $\langle c\rangle\mapsto\langle Tc\rangle$
and on planes by $u^\perp\mapsto (uT^{-1})^\perp$. If $M$ is a union of point orbits $O_1,\dots,O_s$
with multiplicities $x_i$, then the plane intersection numbers are constant on plane orbits. Let
$H_1,\dots,H_t$ be representatives of the plane orbits and $a_{ji}=|O_i\cap H_j|$. Then
\[
  \max\ \sum_{i} |O_i|\,x_i \quad\text{s.t.}\quad \sum_i a_{ji}x_i\le r\ (1\le j\le t),\qquad x_i\in\{0,1\}
\]
yields a point set with all plane intersections $\le r$, i.e.\ an $[n,4,\ge n-r]_q$ code with
$n=\sum|O_i|x_i$. For a target entry $(d,n_{\mathrm{up}}(d))$ of the table, the code is an
improvement iff $n<n_{\mathrm{up}}(n-r)$.

We sampled random $T\in\GL(4,q)$ whose projective order lies in a prescribed set. Orders $8,16,20,24,26,40$
were productive for $q=9$ and $7,14,21$ for $q=8$; these give between $33$ and $108$ point orbits.
We solved the program with HiGHS (single thread, 4--12\,s per instance, not necessarily to
optimality) for every $r$ with an open target. In total about 370 groups were tried. The search
scripts are not part of this package; the resulting codes and the verification scripts are
(Appendix~\ref{app:data}).

\subsection{Greedy lengthening}
If $M$ has maximum plane intersection $r$ and $P$ is a point all of whose planes contain at most
$r-1$ points of $M$, then $M+P$ (multiplicities allowed) still has maximum plane intersection $r$.
This lengthens the code from $[n,4,d]$ to $[n+1,4,d+1]$. Applying this step greedily to the orbit
solutions produced 13 of the 32 base codes (marked ``orbits + lengthening'' in
\cref{tab:base}); some of them have repeated columns.

\subsection{The base codes}
\Cref{tab:base} lists the 32 base codes. Every improved entry of \cref{tab:q9,tab:q8} is a punctured
code of one of them. The base codes were chosen by a greedy set cover among the $230$ verified code files.
The table shows the construction (order of the prescribed cyclic group, when the code is an orbit
union), the maximum column multiplicity, the number of distinct nonzero weights, the gcd of the nonzero
weights, the maximum weight, the number of table entries the code improves (in parentheses: how many of
them meet the lower bound), and the SHA-256 prefix of the data file
\texttt{data/codes/q<q>\_n<n>\_d<d>.txt}.

{\small
\setlength{\tabcolsep}{4pt}
\begin{longtable}{lllrrrrl}
\caption{The 32 base codes (all certified in Lean).}\label{tab:base}\\
\toprule
code & construction & mult & \#wts & gcd & $w_{\max}$ & entries & SHA-256 (prefix)\\
\midrule
\endfirsthead
\caption[]{(continued) The 32 base codes (all certified in Lean).}\\
\toprule
code & construction & mult & \#wts & gcd & $w_{\max}$ & entries & SHA-256 (prefix)\\
\midrule
\endhead
$[163,4,140]_8$ & $|\langle T\rangle|=14$ & 1 & 9 & 2 & 162 & 4 & \texttt{ab8a14270b42d5f4}\\
$[259,4,224]_8$ & $|\langle T\rangle|=14$ & 1 & 6 & 4 & 252 & 4 & \texttt{a8564ba5338ec784}\\
$[385,4,335]_8$ & $|\langle T\rangle|=14$ & 1 & 4 & 1 & 384 & 7 & \texttt{324fd0eaa554cb6b}\\
$[395,4,344]_8$ & $|\langle T\rangle|=14$ & 1 & 3 & 8 & 392 & 8 (8) & \texttt{57f2c76cc0bbb19c}\\
$[397,4,345]_8$ & $|\langle T\rangle|=14$ & 1 & 8 & 1 & 394 & 1 (1) & \texttt{cd6dad28e828a03d}\\
$[407,4,354]_8$ & $|\langle T\rangle|=7$ & 1 & 8 & 2 & 406 & 2 (2) & \texttt{69103d3f0cced90e}\\
$[422,4,368]_8$ & $|\langle T\rangle|=7$ & 1 & 4 & 8 & 400 & 8 (8) & \texttt{fe9075506fa4377f}\\
$[431,4,376]_8$ & $|\langle T\rangle|=7$ & 1 & 4 & 8 & 400 & 8 (8) & \texttt{d8e339c3b2489a27}\\
$[220,4,192]_9$ & $|\langle T\rangle|=20$ & 1 & 5 & 3 & 204 & 3 & \texttt{392a4a3c89c5820e}\\
$[229,4,200]_9$ & orbits + lengthening & 2 & 12 & 1 & 226 & 7 & \texttt{08f00c7af5ce8ec9}\\
$[265,4,232]_9$ & $|\langle T\rangle|=24$ & 1 & 13 & 1 & 260 & 2 & \texttt{b6b70d374a2836cd}\\
$[277,4,243]_9$ & $|\langle T\rangle|=24$ & 1 & 8 & 3 & 273 & 3 & \texttt{caca63a3daca3d84}\\
$[292,4,256]_9$ & orbits + lengthening & 2 & 12 & 1 & 290 & 4 & \texttt{72e035ddbf489fdc}\\
$[303,4,266]_9$ & orbits + lengthening & 2 & 10 & 1 & 301 & 5 & \texttt{13ee0205379bb524}\\
$[314,4,276]_9$ & orbits + lengthening & 2 & 11 & 1 & 312 & 10 & \texttt{d62a7353317e2154}\\
$[318,4,279]_9$ & orbits + lengthening & 3 & 6 & 3 & 315 & 3 & \texttt{ecc272c062828ccf}\\
$[333,4,292]_9$ & $|\langle T\rangle|=24$ & 1 & 16 & 1 & 333 & 4 & \texttt{d7da904ca001abf6}\\
$[345,4,303]_9$ & orbits + lengthening & 2 & 8 & 3 & 342 & 6 & \texttt{93b8f5fd281d85f9}\\
$[355,4,312]_9$ & orbits + lengthening & 3 & 7 & 3 & 351 & 9 & \texttt{c30d4cf4cead8680}\\
$[382,4,336]_9$ & orbits + lengthening & 2 & 7 & 3 & 360 & 4 & \texttt{636b151450aa9b01}\\
$[390,4,343]_9$ & orbits + lengthening & 1 & 13 & 1 & 385 & 7 & \texttt{a735efecc63eee1b}\\
$[408,4,360]_9$ & $|\langle T\rangle|=16$ & 1 & 3 & 9 & 378 & 17 & \texttt{17659a7cc1f6e578}\\
$[422,4,372]_9$ & orbits + lengthening & 2 & 7 & 3 & 420 & 3 & \texttt{f9aae5b910d2ada1}\\
$[426,4,375]_9$ & orbits + lengthening & 1 & 14 & 1 & 425 & 3 & \texttt{7cb0da2406c65270}\\
$[439,4,387]_9$ & orbits + lengthening & 3 & 3 & 9 & 432 & 9 & \texttt{3d90eb9cf624b466}\\
$[468,4,414]_9$ & $|\langle T\rangle|=26$ & 1 & 4 & 9 & 468 & 18 & \texttt{b20bdceedfbef563}\\
$[473,4,417]_9$ & orbits + lengthening & 2 & 7 & 3 & 456 & 3 & \texttt{f62b5bdae030557d}\\
$[486,4,429]_9$ & $|\langle T\rangle|=16$ & 1 & 6 & 3 & 453 & 6 & \texttt{c3d4ea832d31fa01}\\
$[496,4,438]_9$ & $|\langle T\rangle|=8$ & 1 & 9 & 3 & 480 & 6 & \texttt{f7f0d15a44168321}\\
$[553,4,489]_9$ & $|\langle T\rangle|=24$ & 1 & 8 & 1 & 552 & 3 & \texttt{8e9a265ada7c3941}\\
$[565,4,500]_9$ & $|\langle T\rangle|=24$ & 1 & 6 & 1 & 528 & 5 & \texttt{83417062ebf18cb2}\\
$[579,4,513]_9$ & $|\langle T\rangle|=16$ & 1 & 3 & 9 & 576 & 9 (2) & \texttt{4e272d7d4ec1953b}\\
\bottomrule
\end{longtable}

}

\subsection{Three examples}
Field elements are written as integers. In $\F_8=\F_2[x]/(x^3+x+1)$, the integer $c_0+2c_1+4c_2$
denotes $c_0+c_1x+c_2x^2$. In $\F_9=\F_3[x]/(x^2+x+2)$, the integer $c_0+3c_1$ denotes $c_0+c_1x$.
The matrices act on column vectors.

\begin{itemize}
\item $[431,4,376]_8$ is a union of $4$ fixed points and $61$ orbits of size $7$ of
$T=\left(\begin{smallmatrix}3&0&3&7\\7&3&3&1\\0&2&6&5\\6&2&6&6\end{smallmatrix}\right)$ (projective order
$7$). Its weight enumerator is $376^{3577}\,384^{441}\,392^{49}\,400^{28}$ (exponents are the numbers
of codewords), so the code is $8$-divisible.
\item $[395,4,344]_8$ is a union of orbits of sizes $1,2$ and $28\times14$ of
$T=\left(\begin{smallmatrix}2&5&5&6\\2&4&0&0\\6&0&6&7\\4&6&0&0\end{smallmatrix}\right)$ (order $14$). Its
weight enumerator is $344^{3255}\,352^{833}\,392^{7}$, again $8$-divisible.
The code $[422,4,368]_8$ (order $7$) has weights $368,376,392,400$.
\item $[579,4,513]_9$ is a union of orbits of sizes $1,2$ and $36\times 16$ of
$T=\left(\begin{smallmatrix}7&4&5&5\\4&1&8&3\\7&4&5&2\\6&5&6&1\end{smallmatrix}\right)$ (order $16$). Its
weight enumerator is $513^{5336}\,522^{1216}\,576^{8}$, so the code is $9$-divisible.
\end{itemize}
We have no conceptual explanation of these divisible codes. They would be natural starting points for
a geometric (e.g.\ minihyper-based) description.

\section{Verification}\label{sec:lean}

\subsection{Exhaustive enumeration}
For every code file, a short Python program (\texttt{data/verify.py}) enumerates all $q^4-1$ nonzero
messages, computes the codewords with table-based $\F_q$ arithmetic, and prints the minimum distance
and the full weight distribution. All 32 base codes, and all 230 code files of the search, have the
claimed minimum distance. An independent re-implementation with its own $\F_8$/$\F_9$ arithmetic
confirmed all 230 files. We also recomputed the minimum distances directly from the column lists
embedded in the Lean file, in the field representation used there.

\subsection{Lean formalization}
The file \texttt{lean/Main.lean} is self-contained Lean~4.34.1 (core library only, no Mathlib), with
one namespace per field. It is checked by running \texttt{lean Main.lean} in the directory
\texttt{lean/} with the Lean~4.34.1 toolchain (no project file and no Mathlib are needed).

\paragraph{Fields.} $\F_9$ is realised as $\F_3[i]/(i^2+1)$ (the polynomial $x^2+1$ is irreducible over
$\F_3$). The codes were found over $\F_3[x]/(x^2+x+2)$ and were transported by the field isomorphism
$x\mapsto 1+i$, which preserves Hamming weights. $\F_8$ is realised as $\F_2[x]/(x^3+x+1)$ with three
Boolean coefficients.
\begin{lstlisting}
structure GF9 where
  re : Fin 3
  im : Fin 3
instance : Add GF9 := ⟨fun x y => ⟨x.re + y.re, x.im + y.im⟩⟩
instance : Mul GF9 := ⟨fun x y => ⟨x.re * y.re - x.im * y.im, x.re * y.im + x.im * y.re⟩⟩
\end{lstlisting}
The theorems \texttt{Q9.GF9\_is\_field} and \texttt{Q8.GF8\_is\_field} state, and prove by exhaustive
\texttt{decide}, associativity and commutativity of $+$ and $\cdot$, the neutral elements, additive
inverses, distributivity, $1\ne0$, and $\forall x\ne0\ \exists y,\ xy=1$.

\paragraph{The statement.} The predicate used in the main theorems is, verbatim (namespace \texttt{Q9};
\texttt{Q8} is identical with \texttt{GF8}):
\begin{lstlisting}
abbrev Vec4 := GF9 × GF9 × GF9 × GF9
def dot (u c : Vec4) : GF9 := u.1 * c.1 + u.2.1 * c.2.1 + u.2.2.1 * c.2.2.1 + u.2.2.2 * c.2.2.2
def codeword (L : List Vec4) (u : Vec4) : List GF9 := L.map (dot u)
def wt (v : List GF9) : Nat := v.countP (fun x => decide (x ≠ 0))
def IsCode (L : List Vec4) (n d : Nat) : Prop :=
  L.length = n ∧ ∀ u : Vec4, u ≠ (0, 0, 0, 0) → d ≤ wt (codeword L u)
\end{lstlisting}
The main theorems are long conjunctions, one conjunct per improved table entry:
\begin{lstlisting}
theorem Q9.main9 :
    (∃ L : List Vec4, IsCode L 218 190) ∧ (∃ L : List Vec4, IsCode L 219 191) ∧ …
    ∧ (∃ L : List Vec4, IsCode L 579 513)                  -- 149 conjuncts
theorem Q8.main8 :
    (∃ L : List Vec4, IsCode L 160 137) ∧ … ∧ (∃ L : List Vec4, IsCode L 431 376)   -- 42 conjuncts
\end{lstlisting}
The conjuncts are exactly the pairs $(n,d)$ of \cref{tab:q9,tab:q8}.

\paragraph{The certificate and its soundness.} For each base code $C$ the file contains the column list
(encoded as base-$q$ integers and decoded by an explicit function \texttt{ofCode}) and the theorem
\begin{lstlisting}
theorem C_431_376_cert : certify C_431_376 431 376 = true := by decide +kernel
\end{lstlisting}
\texttt{certify L n d} checks \texttt{L.length = n} and, for each of the $(q^4-1)/(q-1)$ normalized
messages (first nonzero coordinate $1$), that the support of the codeword has at least $d$ elements.
The codeword is computed in \emph{bit-sliced} form. A vector of $\F_q^n$ is stored as $2$ bit masks per
$\F_3$-component (for $\F_9$) or $3$ bit masks (for $\F_8$), and field addition and scalar
multiplication become a few \texttt{\&\&\&}, \texttt{|||}, \texttt{\^{}\^{}\^{}} operations on natural
numbers, which the Lean kernel evaluates with GMP arithmetic. The soundness theorem
\begin{lstlisting}
theorem isCode_of_certify (L : List Vec4) (n d : Nat) (h : certify L n d = true) : IsCode L n d
\end{lstlisting}
is proved in Lean as follows.
\begin{enumerate}
\item Bit-level lemmas (\texttt{V3.add\_spec}, \texttt{V3.smul\_spec}, \texttt{mk3\_spec},
\texttt{nz\_spec}) show that bit $c$ of the computed masks encodes $u\cdot c_c$.
\item The lemma \texttt{countP\_eq\_pc} shows that the bit count equals \texttt{wt (codeword L u)}.
\item A projective reduction (\texttt{wt\_smul}, using field identities proved by \texttt{decide}) extends
the bound from normalized messages to all $u\neq0$.
\end{enumerate}
Puncturing (Lemma~\ref{lem:punct}) is the Lean theorem
\begin{lstlisting}
theorem isCode_drop (L : List Vec4) (n d t : Nat) (h : IsCode L n d) :
    IsCode (L.drop t) (n - t) (d - t)
\end{lstlisting}
and each conjunct of \texttt{main9}/\texttt{main8} is proved as
$\langle$\texttt{C.drop t, isCode\_drop \_ \_ \_ t C\_isCode}$\rangle$.

\paragraph{Statement audit.}
\begin{itemize}
\item \texttt{IsCode} quantifies over \emph{all} nonzero $u$. The codeword is the list of the $n$ dot
products, and \texttt{wt} counts nonzero entries.
\item The length is part of the predicate.
\item The coefficient structure is a field with $q$ elements: it is built as \texttt{Fin 3 × Fin 3}, respectively $3$ bits, and the field axioms are proved.
\item The table values $n_{\mathrm{up}}(d)$ are \emph{not} part of the Lean statement. That each pair
$(n,d)$ improves the table is a comparison with the web page (copies kept with the data).
\end{itemize}
The column lists in the Lean file were decoded independently and re-checked by enumeration (see above).

\paragraph{Axioms and resources.} The output of \texttt{\#print axioms}, followed by the resource
summary printed by the wrapper script used for the run (file \texttt{lean/lean\_output.txt}), is:
\begin{lstlisting}
'Q8.GF8_is_field' does not depend on any axioms
'Q8.main8' depends on axioms: [propext, Classical.choice, Quot.sound]
'Q9.GF9_is_field' depends on axioms: [propext]
'Q9.main9' depends on axioms: [propext, Classical.choice, Quot.sound]
leanchk: wall=436.45s maxrss=3865576kB exit=0
\end{lstlisting}
Only the standard axioms appear. There is no \texttt{sorry} and no \texttt{native\_decide}
(no \texttt{Lean.ofReduceBool}): the 32 certificates are evaluated by the kernel
(\texttt{decide +kernel}). The whole file checks in $436$\,s with a peak memory of $3.9$\,GB.
The option \texttt{set\_option Elab.async false} is set only to keep the number of threads, and hence
the address space, small under the checker's memory limit; it has no logical effect.

\section{Limitations and open questions}
\begin{itemize}
\item \emph{Lower bounds are cited, not re-proved.} Corollary~\ref{cor:exact} depends on
\cite[Thm.~1.1(4)]{KanazawaMaruta2011} and \cite[Thm.~1.7(4)]{KOM2017}. Only the upper bounds
(existence of the codes) are formally verified.
\item \emph{Novelty} is asserted relative to Maruta's online table as of October 8, 2026, the
papers~\cite{KanazawaMaruta2011,KOM2017,BFM2021,MaLuo2024,KKM2026}, \cite{Grassl}, \cite{MinT}, and a
literature search. It is ``to our knowledge'' only.
\item The search is heuristic: the integer programs were usually stopped by a time limit, and only
cyclic groups were tried. We did not search $q=7$ or $k\ge5$, nor the short-length entries covered by
Grassl's tables.
\item Most remaining gaps in the mid-range $q=9$ table are still $2$--$3$ wide. For many targets
(e.g.\ $q=9$, $67\le r\le 78$, or $q=8$, $r\in\{48,49,58,59,60\}$) our best codes fall exactly one
point short of the length that would meet the lower bound. It would be interesting to know whether
this is a genuine obstruction.
\item The $8$- and $9$-divisible codes found here (\cref{sec:construction}) suggest an infinite family;
a geometric construction of them is open.
\end{itemize}

{\small
\bibliographystyle{plain}
\bibliography{references}
}

\appendix
\section{All improved table entries}
Columns: $d$; Griesmer bound $g=g_q(4,d)$; the interval $n_{\mathrm{lo}}$--$n_{\mathrm{up}}$ in
Maruta's table~\cite{MarutaTable} as of October 8, 2026; and the new upper bound. Entries marked
$^{*}$ meet the lower bound, so $n_q(4,d)$ is determined (Corollary~\ref{cor:exact}). The base code and
puncturing amount for every row are listed in \texttt{data/improvements.csv}. (For $d=196$, $q=9$, the
lower bound $222$ of the table has meanwhile been raised to $223$ by~\cite[Thm.~1.3(c)]{KKM2026}; the
entry remains open.)

{\small
\setlength{\tabcolsep}{3.5pt}
\begin{longtable}{rrcr|rrcr|rrcr}
\caption{Improved entries of the table of $n_8(4,d)$ (42 entries).}\label{tab:q8}\\
\toprule
$d$ & $g$ & table & new & $d$ & $g$ & table & new & $d$ & $g$ & table & new \\
\midrule
\endfirsthead
\caption[]{(continued) Improved entries of the table of $n_8(4,d)$ (42 entries).}\\
\toprule
$d$ & $g$ & table & new & $d$ & $g$ & table & new & $d$ & $g$ & table & new \\
\midrule
\endhead
137 & 159 & 159--161 & 160 & 335 & 384 & 384--386 & 385 & 363 & 416 & 417--418 & \textbf{417}$^{*}$ \\
138 & 160 & 160--162 & 161 & 337 & 387 & 388--389 & \textbf{388}$^{*}$ & 364 & 417 & 418--419 & \textbf{418}$^{*}$ \\
139 & 161 & 161--163 & 162 & 338 & 388 & 389--390 & \textbf{389}$^{*}$ & 365 & 418 & 419--420 & \textbf{419}$^{*}$ \\
140 & 162 & 162--164 & 163 & 339 & 389 & 390--391 & \textbf{390}$^{*}$ & 366 & 419 & 420--421 & \textbf{420}$^{*}$ \\
221 & 254 & 255--257 & 256 & 340 & 390 & 391--392 & \textbf{391}$^{*}$ & 367 & 420 & 421--422 & \textbf{421}$^{*}$ \\
222 & 255 & 256--258 & 257 & 341 & 391 & 392--393 & \textbf{392}$^{*}$ & 368 & 421 & 422--423 & \textbf{422}$^{*}$ \\
223 & 256 & 257--259 & 258 & 342 & 392 & 393--394 & \textbf{393}$^{*}$ & 369 & 423 & 424--425 & \textbf{424}$^{*}$ \\
224 & 257 & 258--260 & 259 & 343 & 393 & 394--395 & \textbf{394}$^{*}$ & 370 & 424 & 425--426 & \textbf{425}$^{*}$ \\
329 & 378 & 378--380 & 379 & 344 & 394 & 395--396 & \textbf{395}$^{*}$ & 371 & 425 & 426--427 & \textbf{426}$^{*}$ \\
330 & 379 & 379--381 & 380 & 345 & 396 & 397--398 & \textbf{397}$^{*}$ & 372 & 426 & 427--428 & \textbf{427}$^{*}$ \\
331 & 380 & 380--382 & 381 & 353 & 405 & 406--407 & \textbf{406}$^{*}$ & 373 & 427 & 428--429 & \textbf{428}$^{*}$ \\
332 & 381 & 381--383 & 382 & 354 & 406 & 407--408 & \textbf{407}$^{*}$ & 374 & 428 & 429--430 & \textbf{429}$^{*}$ \\
333 & 382 & 382--384 & 383 & 361 & 414 & 415--416 & \textbf{415}$^{*}$ & 375 & 429 & 430--431 & \textbf{430}$^{*}$ \\
334 & 383 & 383--385 & 384 & 362 & 415 & 416--417 & \textbf{416}$^{*}$ & 376 & 430 & 431--432 & \textbf{431}$^{*}$ \\
\bottomrule
\end{longtable}

\begin{longtable}{rrcr|rrcr|rrcr}
\caption{Improved entries of the table of $n_9(4,d)$ (149 entries).}\label{tab:q9}\\
\toprule
$d$ & $g$ & table & new & $d$ & $g$ & table & new & $d$ & $g$ & table & new \\
\midrule
\endfirsthead
\caption[]{(continued) Improved entries of the table of $n_9(4,d)$ (149 entries).}\\
\toprule
$d$ & $g$ & table & new & $d$ & $g$ & table & new & $d$ & $g$ & table & new \\
\midrule
\endhead
190 & 216 & 216--219 & 218 & 307 & 347 & 348--351 & 350 & 398 & 449 & 450--453 & 452 \\
191 & 217 & 217--220 & 219 & 308 & 348 & 349--352 & 351 & 399 & 450 & 451--454 & 453 \\
192 & 218 & 218--221 & 220 & 309 & 349 & 350--353 & 352 & 400 & 451 & 452--455 & 454 \\
194 & 220 & 220--224 & 223 & 310 & 350 & 351--354 & 353 & 401 & 452 & 453--456 & 455 \\
195 & 221 & 221--225 & 224 & 311 & 351 & 352--355 & 354 & 402 & 453 & 454--457 & 456 \\
196 & 222 & 222--226 & 225 & 312 & 352 & 353--356 & 355 & 403 & 454 & 455--458 & 457 \\
197 & 223 & 224--227 & 226 & 333 & 376 & 376--380 & 379 & 404 & 455 & 456--459 & 458 \\
198 & 224 & 225--228 & 227 & 334 & 378 & 378--382 & 380 & 405 & 456 & 457--460 & 459 \\
199 & 226 & 226--229 & 228 & 335 & 379 & 379--383 & 381 & 406 & 459 & 459--462 & 460 \\
200 & 227 & 227--230 & 229 & 336 & 380 & 380--384 & 382 & 407 & 460 & 460--463 & 461 \\
231 & 261 & 262--265 & 264 & 337 & 381 & 381--385 & 384 & 408 & 461 & 461--464 & 462 \\
232 & 262 & 263--266 & 265 & 338 & 382 & 382--386 & 385 & 409 & 462 & 462--465 & 463 \\
241 & 272 & 273--276 & 275 & 339 & 383 & 383--387 & 386 & 410 & 463 & 463--466 & 464 \\
242 & 273 & 274--277 & 276 & 340 & 384 & 384--388 & 387 & 411 & 464 & 464--467 & 465 \\
243 & 274 & 275--278 & 277 & 341 & 385 & 385--389 & 388 & 412 & 465 & 465--468 & 466 \\
253 & 287 & 287--290 & 289 & 342 & 386 & 386--390 & 389 & 413 & 466 & 466--469 & 467 \\
254 & 288 & 288--291 & 290 & 343 & 388 & 388--392 & 390 & 414 & 467 & 467--470 & 468 \\
255 & 289 & 289--292 & 291 & 344 & 389 & 389--393 & 392 & 415 & 469 & 469--472 & 471 \\
256 & 290 & 290--293 & 292 & 345 & 390 & 390--394 & 393 & 416 & 470 & 470--473 & 472 \\
262 & 297 & 297--300 & 299 & 346 & 391 & 391--395 & 394 & 417 & 471 & 471--474 & 473 \\
263 & 298 & 298--301 & 300 & 347 & 392 & 392--396 & 395 & 424 & 479 & 479--482 & 481 \\
264 & 299 & 299--302 & 301 & 348 & 393 & 393--397 & 396 & 425 & 480 & 480--483 & 482 \\
265 & 300 & 300--303 & 302 & 349 & 394 & 394--398 & 397 & 426 & 481 & 481--484 & 483 \\
266 & 301 & 301--304 & 303 & 350 & 395 & 396--399 & 398 & 427 & 482 & 482--485 & 484 \\
267 & 302 & 302--306 & 305 & 351 & 396 & 397--400 & 399 & 428 & 483 & 483--486 & 485 \\
268 & 303 & 303--307 & 306 & 352 & 398 & 398--401 & 400 & 429 & 484 & 484--487 & 486 \\
269 & 304 & 304--308 & 307 & 353 & 399 & 399--402 & 401 & 433 & 489 & 490--492 & 491 \\
270 & 305 & 305--309 & 308 & 354 & 400 & 400--403 & 402 & 434 & 490 & 491--493 & 492 \\
271 & 307 & 307--310 & 309 & 355 & 401 & 401--404 & 403 & 435 & 491 & 492--494 & 493 \\
272 & 308 & 308--311 & 310 & 356 & 402 & 402--405 & 404 & 436 & 492 & 493--495 & 494 \\
273 & 309 & 309--312 & 311 & 357 & 403 & 403--406 & 405 & 437 & 493 & 494--496 & 495 \\
274 & 310 & 310--314 & 312 & 358 & 404 & 405--407 & 406 & 438 & 494 & 495--497 & 496 \\
275 & 311 & 311--315 & 313 & 359 & 405 & 406--408 & 407 & 487 & 550 & 550--552 & 551 \\
276 & 312 & 312--316 & 314 & 360 & 406 & 407--409 & 408 & 488 & 551 & 551--553 & 552 \\
277 & 313 & 313--317 & 316 & 370 & 418 & 418--422 & 420 & 489 & 552 & 552--554 & 553 \\
278 & 314 & 315--318 & 317 & 371 & 419 & 419--423 & 421 & 496 & 560 & 560--562 & 561 \\
279 & 315 & 316--319 & 318 & 372 & 420 & 420--424 & 422 & 497 & 561 & 561--563 & 562 \\
289 & 327 & 328--331 & 330 & 373 & 421 & 421--425 & 424 & 498 & 562 & 562--564 & 563 \\
290 & 328 & 329--332 & 331 & 374 & 422 & 423--426 & 425 & 499 & 563 & 563--565 & 564 \\
291 & 329 & 330--333 & 332 & 375 & 423 & 424--427 & 426 & 500 & 564 & 564--566 & 565 \\
292 & 330 & 331--334 & 333 & 379 & 428 & 429--432 & 431 & 505 & 570 & 570--572 & 571 \\
298 & 337 & 338--341 & 340 & 380 & 429 & 430--433 & 432 & 506 & 571 & 571--573 & 572 \\
299 & 338 & 339--342 & 341 & 381 & 430 & 431--434 & 433 & 507 & 572 & 572--574 & 573 \\
300 & 339 & 340--343 & 342 & 382 & 431 & 432--435 & 434 & 508 & 573 & 573--575 & 574 \\
301 & 340 & 341--344 & 343 & 383 & 432 & 433--436 & 435 & 509 & 574 & 574--576 & 575 \\
302 & 341 & 342--345 & 344 & 384 & 433 & 434--437 & 436 & 510 & 575 & 575--577 & 576 \\
303 & 342 & 343--347 & 345 & 385 & 434 & 435--438 & 437 & 511 & 576 & 576--578 & 577 \\
304 & 343 & 344--348 & 347 & 386 & 435 & 436--439 & 438 & 512 & 577 & 578--579 & \textbf{578}$^{*}$ \\
305 & 344 & 345--349 & 348 & 387 & 436 & 437--440 & 439 & 513 & 578 & 579--580 & \textbf{579}$^{*}$ \\
306 & 345 & 346--350 & 349 & 397 & 448 & 449--452 & 451 &  & & &  \\
\bottomrule
\end{longtable}

}

\section{Data and reproduction}\label{app:data}
\begin{itemize}
\item \texttt{data/codes/}: the 32 base codes. The first line of each file is \texttt{q= k= n= d=}; each
further line \texttt{m x0 x1 x2 x3} is a column, repeated $m$ times, in the field encoding of
\cref{sec:construction}.
\item \texttt{data/improvements.csv}: all 191 entries with base code and puncturing amount.
\item \texttt{data/SHA256SUMS}: hashes of all data and Lean files.
\item \texttt{lean/Main.lean}, \texttt{lean/lean\_output.txt}: the formal certificate and its checker
output.
\end{itemize}
To reproduce the verification, run \texttt{python3 data/verify.py data/codes/q8\_n431\_d376.txt} for
each code (exhaustive enumeration, seconds per code) and \texttt{lean Main.lean} in the directory
\texttt{lean/} with Lean~4.34.1 ($\approx 7$ minutes on a dedicated core, up to about twice that on a
loaded machine; peak memory about $3.9$\,GB).
\end{document}
